
Neural networks trained on spuriously correlated data achieve high average accuracy while failing on minority groups that lack the spurious feature --- a problem typically addressed with group annotations. This work investigates whether the per-sample normalized last-layer gradient norm, $\tilde{g}_i = \|\nabla _W \ell _{i}\| / \|h(x_{i})\|$, constitutes a usable label-free minority group proxy signal during standard ERM training. The hypothesis is that minority samples, unable to exploit the spurious shortcut, maintain persistently elevated prediction residuals that appear directly in their gradient norms via the outer-product decomposition of the fully connected layer gradient. On the Waterbirds benchmark, $\tilde{g}$ achieves AUC = 0.914 for minority group prediction at epoch 5 of ERM training --- without any group annotations --- with a minority/majority ratio of 8.8x that persists to 8.5x at epoch 10. Feature norm distributions (h\_norm\_std\_ratio $\approx$ 0.10) confirm that the signal reflects prediction-error differences rather than feature-scale artifacts. The overall gate result for the existence experiment is PARTIAL (2/3 criteria satisfied): the ratio and AUC criteria pass, while a third criterion --- balance deviation of the top-25\% subset --- fails at 0.379 against a $\leq 0.10$ threshold. Post-hoc analysis identifies this as a criterion design mismatch rather than a mechanism failure. Worst-group accuracy evaluation via a proposed Stage 2 last-layer retraining procedure was not executed and constitutes future work.

